\section*{The results and their mechanisms}
We use the marked family of Paper~I, with its repaired circle quotient
$e:X\to S^3$ and algebraic reduction $f:X\to\PP^1$.
The integral lattice, its monodromies and the positive normal-circle
orientation are fixed in the transport dictionary below.  In particular,
an affine logarithmic parameter is an actual vector in that lattice;
its integral part remains in the boundary gluing map.  The maps in this
paper are specified with these markings.  Current proof checks and the
historical checks are distinguished in the record attached to each full
statement.

\subsection*{1. The repaired circle quotient and the freedom of its extension}
\RungZero{The quotient of the regular family extends across the cusp through
the explicit target solid torus constructed in Paper~I, with the
target and collar conditions fixing the extension used here.
Agreement only on the exterior permits all twelve elements of
$\pi_6(S^3)$, because the relative pinch of an interior six-ball
adds any chosen element (\cref{mapsI:src:A.2.1,maps:src:A.2.2}).}
\Needspace{5\baselineskip}
\RungOne
\begin{enumerate}[label=\arabic*.,leftmargin=*]
\item Use the primitive circle quotient and its explicit period
comparison from Paper~I.  The comparison preserves the signed primitive
covector, so the quotient and its boundary meridian are fixed maps.
\item The cusp target is the solid torus supplied by the local quotient
construction.  Its higher homotopy group in degree six is zero, which
gives the relative uniqueness within this prescribed target.
\item Choose a closed six-ball in the interior of the cusp
neighbourhood.  Collapse its boundary to obtain the relative pinch
$X\to X\vee S^6$, equal to the identity off the ball.
\item Compose the first summand with $e$ and the second with a
representative of $\alpha\in\pi_6(S^3)$.  The resulting map agrees
with $e$ on the exterior and represents $[e]+\alpha$.
This proves the exact scope of the extension assertion.
\end{enumerate}

\subsection*{2. The class of the map to the three-sphere}
\RungZero{The marked quotient has
$[e]=6\nu'=\eta_3\eta_4\eta_5$, the element of order two in
$\pi_6(S^3)=\Z/12$.  A regular inverse image is a three-torus;
the monodromy calculation identifies its normal framing with its Lie
framing, whose real $e$-invariant is $1/2$
(\cref{maps:src:B.1.1,maps:src:B.1.2}).}
\RungOne
\begin{enumerate}[label=\arabic*.,leftmargin=*]
\item Express the framing difference on the regular four-torus.
Its degree-one obstruction is the common invariant covector $\gamma$
modulo two.  Restriction to $\ker\gamma$ therefore lifts this
comparison to $\operatorname{Spin}$.
\item Twice the remaining degree-three obstruction is the restriction
of the common invariant class dual to the invariant vector $d$.
Since $d\in\ker\gamma$, this restriction is zero.  The integral
third cohomology of the three-torus is torsion-free, so the obstruction
itself is zero.
\item Pontryagin--Thom now represents $e$ by the Lie-framed
three-torus.  The periodic flat Dirac operator has paired nonzero
spectrum and a two-dimensional kernel.  The normalization in
\cite[Theorem~4.14, p.~423]{mapsAPS2} gives real $e$-invariant $1/2$.
\item Toda's stabilization calculation identifies the injected
$\pi_6(S^3)$ with the subgroup generated by $2\nu$ in the stable
three-stem.  The class with this invariant is $6\nu'$, also the
displayed composite of the three Hopf maps
\cite[Theorem~4.1 and Corollaries~4.2--4.3]{mapsToda52}.
\end{enumerate}

\subsection*{3. Null-homotopy of the map to the base}
\RungZero{The holomorphic algebraic reduction $f:X\to\PP^1$ is
null-homotopic.  Its explicit factorization through the circle quotient
contains the normalized power map of degree $3\cdot4=12$ on $S^3$,
which kills $\pi_6(S^3)$ (\cref{maps:src:B.2.1}).}
\RungOne
\begin{enumerate}[label=\arabic*.,leftmargin=*]
\item In the two multiplicative coordinates, the normalized power map
is $P_{3,4}(z_1,z_2)=(z_1^3,z_2^4)/\|(z_1^3,z_2^4)\|$.
A regular value has twelve positively oriented inverse images, so its
degree is twelve.
\item The marked quotient construction gives the equality of maps
$f=H\circ P_{3,4}\circ e$, where $H:S^3\to S^2$ is the Hopf
map and $\PP^1$ is identified with $S^2$ by the stated orientation.
\item The degree-twelve self-map acts by multiplication by twelve on
$\pi_6(S^3)=\Z/12$.  Thus $P_{3,4}\circ e$ is null-homotopic;
postcomposition with $H$ gives the required null-homotopy of $f$.
The body also gives the normalized-section homotopy directly.
\end{enumerate}

\subsection*{4. The cusp map retaining the angular coordinate}
\RungZero{The cusp collapse retaining the angular coordinate represents
$\eta_4\eta_5\ne0$ in $\pi_6(S^4)$.  At a triple point its
regular inverse image is a two-torus, and the explicit normal-frame
comparison makes this the Lie-framed torus of Arf invariant one
(\cref{maps:src:B.3.1,maps:src:B.3.2,maps:src:B.3.3}).}
\RungOne
\begin{enumerate}[label=\arabic*.,leftmargin=*]
\item Define the map in the triple-point chart using the retained
angle and the two radial coordinates.  Keep the bounded period
correction when differentiating the normal frame.
\item The resulting frame differs from the product frame by
$\operatorname{diag}(e^{-i(a+b)},e^{ia},e^{ib})$ in
$\operatorname{SU}(3)$.  The displayed two $\operatorname{SU}(2)$
factor homotopies contract this comparison.
\item The inverse-image torus therefore has its Lie framing.  Its
spin quadratic form has Arf invariant one, so its Pontryagin--Thom
class is the nonzero $\eta_4\eta_5$.
\item Two allowed real-period lifts agree on the exterior collar;
their straight-line interpolation gives a relative homotopy.
The core-collapse comparison factors through the specified
$S^3\vee S^2$ cofiber, with the map using its $S^3$ summand.
These are the explicit comparisons among the cusp maps.
\end{enumerate}

\subsection*{5. A finite cell model compatible with the fibration}
\RungZero{The regular torus mapping cylinders, the finite affine caps
and the marked honeycomb cusp fit into a finite CW model with cell
counts $(5,20,45,56,42,17,3)$.  The integral attaching maps give the
homotopy type of the six-sphere
(\cref{maps:src:B.5.1,maps:src:B.5.2}).}
\RungOne
\begin{enumerate}[label=\arabic*.,leftmargin=*]
\item Construct the regular mapping-cylinder bundle using the actual
integral monodromies.  In every finite affine cap retain the drift
from the logarithmic translation; it is part of the attaching map.
\item In the cusp write the seam map as
$[a,b,\theta]\mapsto[a,b-\theta(B_0a-v(a))]$.
The frozen-period comparison is an explicit logarithmic-coordinate
homeomorphism and agrees with the prescribed boundary markings.
\item The positive quotient has interior $T^2\times(0,1)$ and the
two marked boundary tori.  Its incompressible boundary and the
induced fundamental-group isomorphism meet Waldhausen's
boundary-fixed product criterion
\cite[Theorem~6.1, p.~77]{mapsWaldhausen68}.
\item Triangulate these marked pieces and their common faces.  The
resulting mapping-cylinder cells have the displayed counts, and
their alternating sum is two.
\item Van Kampen gives the relations $x^3=x^4=1$.  The primitive
finite pullback lattices and the two cusp sweeps give unimodular
integral attachments, so $H_1=H_2=0$.  Duality, the Euler number,
Hurewicz and Whitehead finish the sphere identification.
\end{enumerate}

\subsection*{6. The local Kummer cover and its global obstruction}
\RungZero{The prescribed degree-$N^2$ relative Kummer cover of the cusp,
for $N>1$, has no extension to a smooth finite cover of the whole
threefold, because its marked monodromy continuation has primitive
ray indices that violate the full-fan smoothness equations.
The explicit semidirect cover and its normal monomial completion
give the finite flat normal completion
(\cref{maps:src:C.1.1,maps:src:C.1.2,maps:src:C.1.3,maps:src:C.1.4,maps:src:C.1.5}).}
\RungOne
\begin{enumerate}[label=\arabic*.,leftmargin=*]
\item On the $N$th root of the base use the actual period
$C(v^N)/N$.  The monomial maps and their lattice inclusions
separate the absolute degree-$N^3$ map from its relative
degree-$N^2$ map.
\item Across adjacent unimodular triangles, smoothness imposes the
displayed divisibility relations on primitive ray indices.
The reverse relations force equality, so a smooth full-fan cover
must use the scalar toric sublattice $n\Z^3$.
\item Continue the prescribed local Kummer lattice by the marked
order-four monodromy.  Its toric intersection is
$\Z w\oplus N\Z d$.  Its primitive ray indices violate that
scalar condition for $N>1$.
\item The explicit word
$A_2^2A_1A_2^2A_1^{-1}$ also produces $N^2$ distinct based
lattices; their residues are exhibited in the body.  The finite
semidirect quotient supplies a global cover with the required
local lattice.
\item Its normal monomial completion has a free basis indexed by
fractional exponents in a half-open fundamental box.  This proves
finite flatness of that normal completion while recording the
smoothness obstruction above.
\end{enumerate}

\subsection*{7. The cyclic universal covers}
\RungZero{For the original marked family with fundamental group
$\Z/d$, where $d$ is prime to six, the universal cover has
$H_2=\Z^{d-1}\oplus\Z/d$, $H_3=\Z/d$, and Euler number
$2d$.  Its cusp divisor products and finite-coefficient cubic
evaluations determine the complete Zhubr invariant, hence one
oriented diffeomorphism type $M_d$ for each $d$
(\cref{maps:src:D.1.1,maps:src:D.1.4,maps:src:D.1.6,maps:src:D.1.9,maps:src:D.1.11}).}
\RungOne
\begin{enumerate}[label=\arabic*.,leftmargin=*]
\item The connected cusp cover has $d$ components in its periodic
normalization.  Its Laurent relation matrix reduces integrally
to $\operatorname{diag}(1,t-1,t-1,0)$.
\item The nontrivial deck characters make the regular torus Koszul
complex acyclic.  The primitive finite-cap degree-two pullbacks
and the cusp sweep with coefficient $d$ determine the integral
groups, including their torsion.
\item The $d$ cusp divisors give classes $x_i$ with
$\sum x_i=0$ and
$x_ix_j=(\mathbf1_{j=i+1}+\mathbf1_{j=i-1}-2\mathbf1_{j=i})\kappa$,
where $\kappa$ generates the degree-four torsion.
\item For the finite-coefficient class $u$ the body computes
$\langle ux_ix_j,[M_d]\rangle$ by the same cyclic matrix,
$\langle u^2x_i,[M_d]\rangle=-12$, and
$\langle u^3,[M_d]\rangle=0$ modulo $d$.  The support-function
restriction has square $-12$, which fixes the second value.
\item Transfer kills the integral pullback in degree four, so
$c_2=p_1=0$.  Multiplication by four is injective on
$\Z^{d-1}\oplus\Z/d$, giving zero refined Pontryagin class.
The full $H_6(K(H_2,2))$ calculation shows that these evaluations
specify the remaining invariant; Zhubr's classification applies
to the already constructed covers.
\item Euler multiplicativity gives $\chi=2d$.  A connected sum of
$g$ copies of $S^3\times S^3$ has Euler number $2-2g$; for
$d>1$ the equality would require $g=1-d<0$.  This supplies the
numerical obstruction to that proposed connected-sum type.
\end{enumerate}

\subsection*{8. The other elliptic surfaces and their six-manifolds}
\RungZero{The primitive specialization lattices and actual normal-circle-torus
maps classify the marked natural-line sphere outputs on all
twenty-one rank-one Mordell--Weil profiles: $59$ patterns on $50$
initial rows, with twelve profiles of unbounded sphere multiplicity,
four of logarithmic order one and five with no sphere output.
Their integral cokernels and full characteristic evaluations also
identify the minimal simply connected nonspherical outputs and
weight-$d$ Engel manifolds, and the free OS49 Gysin lattice retains
the extra torsion outside the coprime sector
(\cref{maps:src:D.2.28,maps:src:D.2.44,maps:src:D.2.45,maps:src:D.2.46,maps:src:D.2.47,maps:src:D.2.48,night:P3F.34}).}
\RungOne
\begin{enumerate}[label=\arabic*.,leftmargin=*]
\item Fix the initial section pair, the component on which the
linearization vanishes, the marked filling and every affine
logarithmic vector.  The profile is the elliptic-fibre configuration;
a pattern is the entire marked initial row, with its stated
arithmetic conditions.
\item Compute the primitive specialization lattices and the actual
normal-circle-torus maps.  Regular-fibre restriction is carried
through the boundary and Gysin maps before the finite-coefficient
classes are used.  Their integral cokernels retain all torsion.
\item The all-order calculations for the original, III-star,
Liu--Rollenske and Wang families give their explicit determinant
or normal-torus matrices.  The additive and multiplicative local
calculations close the remaining profiles; the complete table in
\cref{maps:src:D.2.28} records every pattern.
\item In the simply connected residual cases retain the spin
characteristic class, Zhubr's refined class and the full
fundamental-class image.  The conic surgery and the marked
$Z_{c,\epsilon}$ surgeries give explicit representatives of the
types in \cref{maps:src:D.2.44,maps:src:D.2.45}.
\item For a period sublattice of the cusp, intersect the component
restrictions on each degree-six del Pezzo normalization.  Keep
loops, repeated edges and repeated vertices in this calculation;
the complete trilinear expression is
\cref{maps:src:D.2.47}.
\item In the simply connected weight-$d$ sector the Seifert
determinant is $P=\pm1$.  The full invariant includes the
Pontryagin square modulo $d\gcd(d,2)$ and the third power
modulo $d\gcd(d,3)$, with evaluations $-12$ and $0$.
Together with the cyclic divisor matrix and refined class zero,
these give the single type $M_d$ in \cref{maps:src:D.2.48}.
\item For a free OS49 cap $(m,n,d)=(3k,n,d)$ with $n\ne0$,
put $g=\gcd(3k,n,3d)$.  The exact primitive Gysin lattice has
Smith invariants $g,|n|$, giving
$H_2=H_3=\Z/|n|\oplus\Z/g$ with $g\in\{1,3\}$.
For example $(9,3,1)$ gives $(\Z/3)^2$.
The zero-$n$ calculation and the precise freeness conditions are
retained in \cref{night:P3F.33,night:P3F.34}.
\end{enumerate}

\subsection*{9. The conifold transition}
\RungZero{Moving the zero of the line-bundle linearization gives the local
equation $xy=u(zv-s)$ and a proper analytic family with one ordinary
double point. The actual replacement of $S^2\times D^4$ by
$D^3\times S^3$ and its primitive boundary maps identify its two
smooth sides as the original six-sphere and $S^3\times S^3$,
the latter with algebraic dimension one
(\cref{maps:src:D.3.1,maps:src:D.3.2,maps:src:D.3.3}).}
\RungOne
\begin{enumerate}[label=\arabic*.,leftmargin=*]
\item Start with the marked elliptic surface and the line
$M=\OO(P-O)$.  Move the linearization as $\phi_s=\phi_0-s\phi_1$,
using its exact annular identification on the common exterior.
\item Near the collision the equation is
$xy=u(zv-s)$.  At $s=0$ this is the ordinary double point;
its small resolution has exceptional normal bundle
$\OO(-1)\oplus\OO(-1)$.  The root caps remain in the same
proper family because the collars were chosen uniformly.
\item On real manifolds the transition replaces $S^2\times D^4$
by $D^3\times S^3$.  General position makes the complement
simply connected; Alexander duality and the two primitive boundary
maps give $H_3=\Z^2$ and $H_2=H_4=0$.
\item The spin Wall invariant has the standard rank-two skew form
and zero normal invariant.  The classification, together with
$\Theta_6=0$, identifies the smoothing with $S^3\times S^3$.
\item The root canonical formula and the varying periods give
algebraic dimension one.  The formal cohomology of the exceptional
curve transports the stated structure-sheaf cohomology across the
small resolution.  The comparison with the published CDP quotient
uses its actual algebraic-reduction fibres in
\cref{maps:src:D.3.4}.
\end{enumerate}

\subsection*{10. A fourfold obtained by a transverse fibre product}
\RungZero{The explicit crepant small resolution of the transverse elliptic
fibre product gives a compact complex fourfold with
$K=\Pi^*L^{-5}$ and algebraic dimension two; its Todd and
Pontryagin calculations give signature two and middle form
$\langle1\rangle\oplus\langle1\rangle$.
The computed values $b_1=0$ and $h^{0,1}=1$, preserved by a
K\"ahler modification, obstruct membership in Fujiki's class
$\mathcal C$ through the K\"ahler equality $b_1=2h^{0,1}$
(\cref{maps:src:E.1.2,maps:src:E.2.1}).}
\RungOne
\begin{enumerate}[label=\arabic*.,leftmargin=*]
\item Pull back the Weierstrass family to the root curve of the
threefold fibration.  The integral matrices $A,B$ give the actual
period comparison $AZ_3B^{-1}=\operatorname{diag}(Z_2,s+c_E)$.
\item The unimodular Kuhn fan resolves the local fibre product
without adding rays.  Its period action is proper and its marked
section extends face by face.  The resolution is therefore small
and crepant.
\item The canonical formula on the root curve gives
$K=\Pi^*L^{-5}$.  The split relative elliptic complex and Leray
give $(h^{0,q})=(1,1,0,0,0)$; the explicit meromorphic product
comparison gives algebraic dimension two.
\item A K\"ahler modification in class $\mathcal C$ preserves
$b_1$ and $h^{0,1}$.  Here $b_1=0$ and $h^{0,1}=1$, contradicting
the K\"ahler equality $b_1=2h^{0,1}$.  This proves the stated
exclusion through these invariants.
\item The universal Todd and Pontryagin formulas give
$c_2^2=2$, $c_4=6$ and signature two.  The middle lattice is
positive definite unimodular of rank two, hence has the displayed
form.  The body records the precise scope of each cited formula.
\end{enumerate}

\paragraph{Supporting constructions and remaining scope.}
The later sections keep the complete transverse-pullback,
covering-space and cusp-carrier arguments, including their explicit
quotient maps, and the marked realization of the two collision
profiles.  The remaining classification question is recorded as
\cref{maps:src:D.2.4}; it is kept separate from the proved sectors
and from the current proof-review records.
